\section{有限元空间的性质}\label{sec:finite-element-space-properties}

有限元空间 $\mathcal P_r\Lambda^k(\mathcal T_h)$ 和 $\mathcal P_r^-\Lambda^k(\mathcal T_h)$ 的逼近性质 (当网格尺寸 $h$ 趋于 $0$ 时的渐近性质), 可像通常的有限元分析那样, 利用典范投影局部保持多项式这一事实, 应用 Bramble–Hilbert 引理和尺度变换来建立. [\BookNumber{11}, Theorem 5.3] 详细进行了这一分析, 并证明了下述结果.

\begin{Theorem}\label{thm:canonical-projection-approximation}
以 $\Pi_h$ 表示从 $\Lambda^k(\Omega)$ 到 $\mathcal P_r\Lambda^k(\mathcal T_h)$ 或 $\mathcal P_{r+1}^-\Lambda^k(\mathcal T_h)$ 的典范投影. 设 $1\leq p\leq\infty$ 且 $(n-k)/p<s\leq r+1$. 则 $\Pi_h$ 可有界地延拓至 $W_p^s\Lambda^k(\Omega)$, 且存在不依赖于 $h$ 的常数 $C$, 使得
\[
\|\omega-\Pi_h\omega\|_{L_p\Lambda^k(\Omega)}\leq Ch^s|\omega|_{W_p^s\Lambda^k(\Omega)},\qquad\omega\in W_p^s\Lambda^k(\Omega).
\]
\end{Theorem}

注意, 限制 $s>(n-k)/p$ 来自 Sobolev 嵌入定理: 这是 $W_p^s$ 中的函数在 $k$ 维面上具有迹所需的条件.

对光滑 $k$-形式 $\omega$ 和完全多项式空间 $\mathcal P_r\Lambda^k(\mathcal T_h)$, 可取 $p=2$, $s=r+1$, 得到 $r+1$ 阶的 $L^2$ 估计:
\[
\|\omega-\Pi_{r,h}^k\omega\|_{L^2\Lambda^k}\leq Ch^{r+1}|\omega|_{H^{r+1}\Lambda^k}.
\]
对修剪多项式空间 $\mathcal P_r^-\Lambda^k(\mathcal T_h)$, 它包含完全多项式空间 $\mathcal P_{r-1}\Lambda^k(\mathcal T_h)$, \newline 但不包含 $\mathcal P_r\Lambda^k(\mathcal T_h)$, 阶数降为 $r$:
\[
\|\omega-\Pi_{r,h}^{-,k}\omega\|_{L^2\Lambda^k}\leq Ch^r|\omega|_{H^r\Lambda^k}.
\]

对完全多项式空间, 有上链关系
\[
d\Pi_{r,h}^k\omega=\Pi_{r-1,h}^{k+1}d\omega,
\]
因而
\[
\|d(\omega-\Pi_{r,h}^k\omega)\|_{L^2\Lambda^{k+1}}=\|d\omega-\Pi_{r-1,h}^{k+1}d\omega\|_{L^2\Lambda^{k+1}}\leq Ch^r|d\omega|_{H^r\Lambda^{k+1}},
\]
相对于 $\omega$, $d\omega$ 的收敛阶损失了一阶. 对修剪空间, 上链关系为
\[
d\Pi_{r,h}^{-,k}\omega=\Pi_{r,h}^{-,k+1}d\omega,
\]
因此不损失阶数:
\[
\|d(\omega-\Pi_{r,h}^{-,k}\omega)\|_{L^2\Lambda^{k+1}}\leq Ch^r|d\omega|_{H^r\Lambda^{k+1}}.
\]

现在可以选择一对 Galerkin 子空间 $V_h^{k-1}$, $V_h^k$. 两种显然的选择是
\[
\mathcal P_r\Lambda^{k-1}(\mathcal T_h),\mathcal P_{r-1}\Lambda^k(\mathcal T_h)\qquad\text{and}\qquad\mathcal P_r^-\Lambda^{k-1}(\mathcal T_h),\mathcal P_r^-\Lambda^k(\mathcal T_h),
\]
即两个子空间均用完全多项式空间, 或均用修剪空间. 还有另外两种选择, 分别使用一个完全空间和一个修剪空间. 可将完全多项式空间对中的 $V_h^{k-1}$ 缩小, 得到
\[
\mathcal P_r^-\Lambda^{k-1}(\mathcal T_h),\mathcal P_{r-1}\Lambda^k(\mathcal T_h).
\]
这不改变像 $dV_h^{k-1}$, 且典范投影仍然交换: $d\Pi_{r,h}^{-,k}=\Pi_{r-1,h}^kd$. 类似地, 可将修剪多项式空间对中的 $V_h^{k-1}$ 扩大, 得到
\[
\mathcal P_r\Lambda^{k-1}(\mathcal T_h),\mathcal P_r^-\Lambda^k(\mathcal T_h),
\]
这同样不影响 $d$ 的像. 总之, 对每个次数 $r$, 得到四种可能的选择:
\begin{equation}\label{eq:four-galerkin-form-pairs}
V_h^{k-1}=\left\{\begin{gathered}\mathcal P_r\Lambda^{k-1}(\mathcal T_h)\\\text{or}\\\mathcal P_r^-\Lambda^{k-1}(\mathcal T_h)\end{gathered}\right\},\qquad
V_h^k=\left\{\begin{gathered}\mathcal P_r^-\Lambda^k(\mathcal T_h)\\\text{or}\\\mathcal P_{r-1}\Lambda^k(\mathcal T_h)\ (\text{if }r>1)\end{gathered}\right\}.
\end{equation}

这些选择对变量 $u\in\Lambda^k$ 和 $\sigma=\delta u\in\Lambda^{k-1}$ 所达到的 $L^2$ 逼近阶见 Table~\ref{tab:galerkin-form-approximation-rates}. 由 Theorem~\ref{thm:mixed-improved-errors} 的改进误差估计, 可推出这些估计不仅对 $\sigma$ 和 $u$ 的投影成立, 而且对 Galerkin 解也成立 (当真解光滑时).

当然, 要得到这些结果, 需要满足 Theorem~\ref{thm:mixed-improved-errors} 的假设. 为此必须指定第三个空间 $V_h^{k+1}\subset H\Lambda^{k+1}$. (Galerkin 方法不使用这个空间, 但理论需要它.) 为此取
\[
V_h^{k+1}=\begin{cases}\mathcal P_r^-\Lambda^{k+1}(\mathcal T_h)&\text{if }V_h^k=\mathcal P_r^-\Lambda^k(\mathcal T_h),\\\mathcal P_{r-1}^-\Lambda^{k+1}(\mathcal T_h)&\text{if }V_h^k=\mathcal P_{r-1}\Lambda^k(\mathcal T_h).
\end{cases}
\]

\begin{table}[htbp]\centering
\caption{稳定的 Galerkin 子空间选择所对应的逼近阶.}\label{tab:galerkin-form-approximation-rates}
\begin{tabular}{cc|cccc}\hline
$V_h^{k-1}$&$V_h^k$&$\sigma$&$d\sigma$&$u$&$du$\\\hline
$\mathcal P_r\Lambda^{k-1}$&$\mathcal P_r^-\Lambda^k$&$r+1$&$r$&$r$&$r$\\
$\mathcal P_r^-\Lambda^{k-1}$&$\mathcal P_r^-\Lambda^k$&$r$&$r$&$r$&$r$\\
$\mathcal P_r\Lambda^{k-1}$&$\mathcal P_{r-1}\Lambda^k$&$r+1$&$r$&$r$&$r-1$\\
$\mathcal P_r^-\Lambda^{k-1}$&$\mathcal P_{r-1}\Lambda^k$&$r$&$r$&$r$&$r-1$\\\hline
\end{tabular}
\end{table}

这就定义了 Theorem~\ref{thm:mixed-improved-errors} 所要求的, 相应 de Rham 复形片段的子复形 \eqref{eq:galerkin-subcomplex}. 剩下的假设是存在 $W$-有界上链投影, 即对 $j=k-1$, $k$, $k+1$, 存在有界投影算子 $\pi_h^j:L^2\Lambda^j\to V_h^j$, 在其定义域上与 $d$ 交换. 前面已经指出, 由自由度定义的典范投影构成上链投影, 即在充分光滑的形式上与 $d$ 交换. 然而, 典范投影一般在 $L^2\Lambda^k$ 上无界, 甚至在 $H\Lambda^k$ 上也无界. 例如, $\Pi_{1,h}^0$ 是到分片线性函数空间的通常插值算子. 在超过一维时, 它在 $H\Lambda^0=H^1$ 上无界, 因为它依赖顶点处的点值; 即使在一维时, 它也不是 $L^2$ 有界的. 不过, 有多种方法可构造有界上链投影. [\BookNumber{11}, Section 5.4] 详细介绍的一种方法, 是先施加磨光算子, 再施加典范投影 (遵循 \BookCite{28}), 然后修正所得结果, 得到具有所需性质的投影 (遵循 \BookCite{66}). 这一构造给出了所需的 $W$-有界上链投影, 因而 Theorem~\ref{thm:mixed-basic-error} 和 Theorem~\ref{thm:mixed-improved-errors} 的误差估计成立.

需要指出, 对上链投影还可能希望有其他各种性质, 而刚才回顾的构造并不提供这些性质. 例如, 可能希望该构造是局部的, 或要求界关于次数 $r$ 一致. 因此, 有界上链投影的构造仍是一个活跃的研究领域 \BookCite{30}, \BookCite{43}.
